\documentclass[11pt]{article}
\usepackage[margin=2.5cm]{geometry}
\usepackage{amsmath,amssymb}
\usepackage{enumitem}
\usepackage{hyperref}

\title{Per-Grain Analysis Pipeline: Clustering, Parent-Grain Reconstruction,\\
and Martensite Twin Correspondence}
\author{}
\date{}

\begin{document}
\maketitle

For each region of interest (ROI), corresponding to one reconstructed
austenite grain, the following pipeline is applied. ROIs are selected
manually, by drawing a polygon around what is judged, from the
microstructure map, to correspond to a single pre-deformation austenite
grain; this selection is not an objective, automated segmentation and is
therefore subject to operator error -- an ROI may in reality span more
than one original grain, or capture only part of one, and this
uncertainty propagates into every subsequent step of the analysis.

\section{Misorientation-based sub-grain clustering}

Pixels within the ROI are grouped into sub-grains (``clusters'') using a
misorientation-based region-growing algorithm: neighbouring pixels are
merged when their disorientation is below a threshold angle
\texttt{ang\_thr}, subject to a maximum cluster spatial diameter
\texttt{dmax} and a minimum cluster size \texttt{minidxs}. Each resulting
cluster's average orientation is subsequently \emph{symmetry-branch
consistent} across the whole ROI (\texttt{symmetrize\_clusters}): since a
measured orientation matrix is only defined up to the crystal's point-group
symmetry, all clusters' representative orientations are resolved onto a
single, mutually consistent symmetric branch before any further analysis.

\section{Local misorientation and dislocation-density fields}

Pixel-to-pixel boundary misorientations are computed within the ROI. For
each pixel, the kernel average misorientation $\langle\theta(x)\rangle$ --
the local disorientation averaged over the peripheral pixels at a fixed
kernel radius $x$ -- is evaluated for a range of kernel radii (1st to
$n$th neighbour shell), following the method of
Kamaya~\cite{kamaya2011}, as applied to EBSD data by Moussa et
al.~\cite{moussa2017}. Disorientation values above a threshold
(\texttt{maxmis}) are excluded from the average to prevent grain
boundaries from contaminating the local, sub-grain-scale signal. Assuming
the disorientation gradient is locally constant, $\langle\theta(x)\rangle$
varies linearly with $x$; a linear regression over the sampled kernel
radii yields both the disorientation gradient $\mathrm{d}\langle
\theta\rangle/\mathrm{d}x$ (the slope) and, by extrapolation to $x=0$, an
estimate of the local EBSD measurement noise (the intercept), which would
otherwise be conflated with genuine lattice curvature if the GND density
were instead estimated directly from $\langle\theta\rangle/x$ at a single
kernel radius. The GND density is then obtained from the slope,
\begin{equation}
  \rho_{\mathrm{GND}} = \frac{\alpha}{b}\,\frac{\mathrm{d}\langle\theta\rangle}{\mathrm{d}x},
\end{equation}
with $b$ the Burgers vector magnitude and $\alpha$ a constant depending on
the dislocation character (tilt or twist), evaluated assuming a single
slip system, $\langle100\rangle\{110\}$, in the B2 austenite phase.

\section{Homogeneous lattice direction and provisional zone axis}

The largest cluster in the ROI (by pixel count) is used as an anchor. For a
small set of candidate low-index directions
$\{\langle100\rangle,\langle110\rangle,\langle111\rangle\}$, each
direction's full crystallographically-equivalent family is tested for how
consistently it points along a common sample-frame direction across
\emph{every} cluster in the ROI (not only the anchor): for each symmetric
variant, the direction is expressed in every cluster's own crystal frame
and its angular spread across the pooled cluster set is measured. The
variant with the smallest spread is retained as the ROI's provisional
common lattice direction (``zone axis''), evidence that the sub-grains
share a crystallographic relationship inherited from a common parent.

\section{Twin and kink relationship identification}

Every pair of clusters in the ROI is tested against a library of tabulated
NiTi austenite twinning modes ($\{112\}$--$\{115\}$): for each mode, the
full space of symmetry-equivalent branch choices is searched
(\emph{T-loop}) together with the mode's tabulated twin elements
($K_1$/$\eta_1$/$K_2$), and a pair is accepted as a twin match if the
axis and angle of at least one of the $K_1$-180$^\circ$,
$\eta_1$-180$^\circ$, or $K_2$-shear conditions falls within a tolerance
\texttt{tol\_deg} of the target. Independently, pairs are also tested for
a \emph{kink} relationship: a blind search over low-index candidate axes
(up to a maximum Miller index) for any axis about which the pair's
misorientation is consistent, with no assumed target angle.

\section{Directed parent-orientation reconstruction}

The identified twin and kink relationships define a graph over the ROI's
clusters. Twin edges are additionally \emph{directed}: for each twin pair,
the Schmid factor of the tabulated twin system is evaluated on both sides
under the given loading direction \texttt{axial\_dir}; the side with the
higher (more mechanically favourable) Schmid factor is treated as the
parent-like side, propagation running only parent$\to$child. Kink edges,
for which no independent Schmid factor is defined, remain undirected. A
breadth-first traversal from the best-supported source node (a cluster
with no incoming directed twin edge, reachable to the largest pixel-weighted
subtree) propagates every reachable cluster's orientation back to a common
reference frame -- using the exact tabulated twin operation for twin edges
(denoised relative to the noisy measured relative orientation) and the raw
measured relative orientation for kink-only edges -- and the reconstructed
parent orientation is the pixel-weighted quaternion mean of all these
independent votes.

Because each edge's direction is decided independently, pairwise, from the
two clusters' own Schmid factors, the resulting directed graph is not
guaranteed to be acyclic: three or more mutually twin-related clusters can
in principle form a directed cycle (each locally the more favourable side
relative to its own neighbour, with no cluster left as an unambiguous
overall source). When this occurs, no node qualifies as a clean starting
point for the reconstruction; rather than abort the analysis, the node
with the \emph{fewest} incoming directed twin edges is used as a
best-effort root instead (tie-broken by the largest reachable pixel
weight), and the affected grain is flagged accordingly. This fallback
indicates that the grain's twin network does not reduce to a single
consistent parent-child hierarchy under the Schmid-factor criterion, and
the corresponding reconstruction should be inspected with additional care
-- for example via the per-cluster deviation from the reconstructed
consensus orientation -- as it may indicate that the selected ROI spans
more than one original grain, or that the loading-direction assumption is
not locally applicable.

\section{Verification against tabulated twin geometry}

The reconstructed parent orientation is checked directly against every
cluster in the ROI (independently of the graph used to derive it): treating
the parent as a hypothetical reference orientation, each real cluster is
tested for a twin relationship to it, using the same T-loop/tabulated-mode
search as above, with a tolerance set by the recipe's grain-reconstruction
threshold. This both confirms the internal consistency of the
reconstruction and yields, for every matched cluster, the Schmid factor of
its twin system evaluated on the parent side and on the cluster's own side.

\section{Matrix identification by orientation merging}

Clusters whose true minimum disorientation (searched over the full crystal
symmetry group) to the reconstructed parent orientation falls below the
same threshold are merged into a single ``matrix'' cluster, representing
the (largely) untwinned parent austenite. The merge additionally records
each merged fragment's original disorientation to the reference, for later
inspection.

\section{Consolidated reconstruction report}

For the merged clustering result, every remaining (non-matrix) cluster is
re-classified relative to the confirmed matrix orientation: twin
relationships (with per-mode breakdown, pixel/cluster counts, and Schmid
factors) and, for clusters showing no twin relationship, an additional kink
check. This yields a single consolidated report of how the ROI's total
pixel population partitions into matrix, twin-related, kink-related, and
unidentified fractions.

\section{Martensite correspondence and variant selection}

If the matrix cluster used for merging exists, it serves as the anchor for
a second homogeneous-direction search (restricted to matrix, twin- and
kink-matched clusters only), refining the common $\langle110\rangle$-type
lattice direction shared across the confirmed network; otherwise, the
cluster closest in disorientation to the reconstructed parent is used as a
fallback anchor. This refined direction is mapped, via the tabulated
austenite$\to$martensite direction-correspondence matrices, into
martensite Miller indices for each of the twelve correspondence variants;
variants whose predicted martensite direction is not parallel to the
expected martensite zone axis are discarded. Among the remaining candidate
variants, the one giving the maximum (for tensile loading) transformation
strain along \texttt{axial\_dir} is selected as the most probable
martensite variant retained prior to reverse transformation. Finally, every
confirmed austenite twin in the report is checked for whether its
$K_1$ plane, mapped through the selected variant's correspondence matrix,
is consistent with a known, named martensite twinning plane -- i.e.
whether the austenite twin is itself inherited from a specific martensite
twin system.

\section{Checkpointing}

The clustering result, the merged (matrix-identified) clustering result,
the reconstruction report, the martensite-correspondence result, the
common lattice direction, the anchor cluster, the loading direction, and
the driving recipe are written to a per-grain HDF5 checkpoint, allowing the
full downstream analysis and figure generation to be repeated without
re-running this pipeline.

\begin{thebibliography}{9}

\bibitem{kamaya2011}
M.~Kamaya,
\emph{Assessment of local deformation using EBSD: Quantification of accuracy of measurement and definition of local gradient},
Ultramicroscopy \textbf{111} (2011) 1189--1199.

\bibitem{moussa2017}
C.~Moussa, M.~Bernacki, R.~Besnard, N.~Bozzolo,
\emph{Statistical analysis of dislocations and dislocation boundaries from EBSD data},
Ultramicroscopy \textbf{179} (2017) 63--72.

\end{thebibliography}

\end{document}
